Urban mobility platforms in data-scarce settings must integrate heterogeneous
real-world feeds while preserving transparent provenance and avoiding
synthetic fallbacks. This paper presents \textit{IRIDIUM}, a
provenance-aware urban mobility platform for Azerbaijani cities, with Baku
as the baseline deployment case. IRIDIUM constructs a request-time digital
twin of the transport network from configured sources, attaches per-source
provenance to every response, and exposes a machine-readable
\texttt{data\_status} field indicating whether returned content is
\texttt{live}, \texttt{recorded}, \texttt{config\_required},
\texttt{permission\_required}, or \texttt{unavailable}. On top of this
twin, the platform defines modular services for congestion forecasting,
multi-objective routing, district-level mobility equity scoring, anomaly
detection, and a federated-learning extension path for privacy-preserving
model training.

The paper formalizes the mobility graph, routing objective,
spatio-temporal forecasting target, equity score, anomaly severity
function, and FedAvg aggregation rule. The current implementation
integrates OpenStreetMap through a PostgreSQL~16/PostGIS~3.4 backend and
ingests live Open-Meteo observations for Baku and Quba; traffic and transit
adapters are registered but remain inactive in the baseline configuration.
Evaluation focuses on three verifiable aspects: API semantics, end-to-end
response latency, and real weather integration. Across $n=50$ loopback
requests per endpoint, median latencies were \SI{2}{ms} for
\texttt{/health}, \SI{45}{ms} for \texttt{/network/graph}, and
\SI{55}{ms} for \texttt{/forecast/congestion}. The weather pipeline
captured a mean Baku--Quba temperature differential of
$\overline{\Delta T} = 2.50 \pm 0.82$\,\textdegree{}C and approximately
16\,mm cumulative precipitation over 72 hours in Baku. Forecasting
accuracy and multimodal-routing benchmarks are intentionally deferred until
traffic histories and transit-feed agreements become available. These
results position IRIDIUM as a transparent baseline architecture for urban
mobility analytics under partial data availability.
